\documentclass[12pt]{article}
\usepackage{amsmath, amssymb}
\usepackage[margin=1in]{geometry}
\setlength{\parskip}{6pt}

\title{QUBO Formulation: Hamiltonian Cycle Problem}
\date{}

\begin{document}
\maketitle

\subsection*{Hamiltonian Cycle Problem}

\paragraph{Problem Statement.}
Given an undirected graph $G=(V,E)$ with $N=|V|$ vertices, find a cycle that visits every vertex exactly once and returns to the start. Equivalently, select a subset of edges such that every vertex has degree exactly two and the selected edges form a single connected component. The objective is to minimize the number of selected edges that are not present in the input graph, with a value of zero indicating the existence of a valid Hamiltonian cycle.

\paragraph{Decision Variables.}
For each pair of distinct vertices $i,j \in V$ with $i<j$, introduce a binary variable $x_{ij} \in \{0,1\}$, where $x_{ij}=1$ if edge $(i,j)$ is included in the cycle and $0$ otherwise. To convert the subtour elimination inequalities into equalities suitable for a QUBO penalty formulation, introduce a non-negative integer slack variable $s_S$ for each proper non-empty subset $S \subset V$. In a binary encoding, $s_S$ may be represented by one or more binary variables depending on the maximum possible deficit $|S|-1$.

\paragraph{QUBO Derivation.}
\textbf{Step 1 --- Objective Function.}
The original objective minimizes the number of selected edges that do not belong to $E$:
\begin{equation}
\min_{x} \sum_{(i,j) \notin E} x_{ij}.
\end{equation}

\textbf{Step 2 --- Constraints.}
The problem is subject to two families of constraints:
\begin{align}
\sum_{j \in V, j \neq i} x_{ij} &= 2, \quad \forall i \in V, \label{eq:degree} \\
\sum_{(i,j) \in E(S)} x_{ij} + s_S &= |S| - 1, \quad \forall S \subset V, S \neq \emptyset, V, \label{eq:subtour}
\end{align}
where $E(S) = \{(i,j) \in E : i,j \in S, i<j\}$ denotes the set of candidate edges within subset $S$.

\textbf{Step 3 --- Penalty Term Expansion.}
We convert each constraint into a quadratic penalty term using positive weights $\lambda$ and $\mu$.

\textit{(a) Degree constraint penalty:}
\begin{equation}
P_{\text{deg}} = \lambda \sum_{i \in V} \left( \sum_{j \in V, j \neq i} x_{ij} - 2 \right)^2.
\end{equation}
\textit{(b) Algebraic expansion:}
Expanding the square and applying the idempotent property $x_{ij}^2 = x_{ij}$ for binary variables yields:
\begin{align}
\left( \sum_{j \neq i} x_{ij} - 2 \right)^2 &= \sum_{j \neq i} x_{ij}^2 - 4 \sum_{j \neq i} x_{ij} + 4 + \sum_{j \neq k, j,k \neq i} x_{ij} x_{ik} \nonumber \\
&= -3 \sum_{j \neq i} x_{ij} + 4 + \sum_{j \neq k, j,k \neq i} x_{ij} x_{ik}.
\end{align}
\textit{(c) General contribution:}
Substituting back, the degree penalty becomes:
\begin{equation}
P_{\text{deg}} = \lambda \sum_{i \in V} \left( -3 \sum_{j \neq i} x_{ij} + 4 + \sum_{j \neq k, j,k \neq i} x_{ij} x_{ik} \right).
\end{equation}

\textit{(a) Subtour constraint penalty:}
Let $k_S = |S|-1$. The penalty is:
\begin{equation}
P_{\text{sub}} = \mu \sum_{S \subset V} \left( \sum_{(i,j) \in E(S)} x_{ij} + s_S - k_S \right)^2.
\end{equation}
\textit{(b) Algebraic expansion:}
Expanding the square and using $s_S^2 = s_S$ (for binary slack encoding) gives:
\begin{align}
\left( \sum_{(i,j) \in E(S)} x_{ij} + s_S - k_S \right)^2 &= \sum_{(i,j) \in E(S)} x_{ij} + s_S - 2k_S \sum_{(i,j) \in E(S)} x_{ij} - 2k_S s_S + k_S^2 \nonumber \\
&\quad + 2 \sum_{(i,j)<(p,q) \in E(S)} x_{ij} x_{pq} + 2 s_S \sum_{(i,j) \in E(S)} x_{ij}.
\end{align}
\textit{(c) General contribution:}
Grouping linear and quadratic terms:
\begin{equation}
P_{\text{sub}} = \mu \sum_{S \subset V} \left[ (1-2k_S) \sum_{(i,j) \in E(S)} x_{ij} + (1-2k_S) s_S + k_S^2 + 2 \sum_{(i,j)<(p,q) \in E(S)} x_{ij} x_{pq} + 2 s_S \sum_{(i,j) \in E(S)} x_{ij} \right].
\end{equation}

\textbf{Step 4 --- Combined Hamiltonian.}
The full QUBO Hamiltonian $H(x,s)$ is the sum of the objective and penalty terms:
\begin{equation}
H(x,s) = \sum_{(i,j) \notin E} x_{ij} + P_{\text{deg}} + P_{\text{sub}}.
\end{equation}

\textbf{Step 5 --- Q Matrix Entries and Offset.}
Reading off the coefficients from the expanded form, the QUBO matrix entries $Q_{u,v}$ and constant offset $c$ are given in closed form as:
\begin{align}
Q_{ij,ij} &= \mathbb{1}((i,j) \notin E) - 6\lambda + \mu \sum_{S \subset V, i,j \in S} (1 - 2(|S|-1)), \\
Q_{ij,kl} &= \lambda \sum_{v \in V} \mathbb{1}(\{i,j\} \cap \{k,l\} = \emptyset \land i,j,k,l \neq v) + 2\mu \sum_{S \subset V} \mathbb{1}(\{i,j\},\{k,l\} \in E(S)), \\
Q_{s_S, s_S} &= \mu (1 - 2(|S|-1)), \\
Q_{ij, s_S} &= 2\mu \cdot \mathbb{1}((i,j) \in E(S)), \\
c &= 4\lambda N + \mu \sum_{S \subset V} (|S|-1)^2,
\end{align}
where $\mathbb{1}(\cdot)$ denotes the indicator function, and all sums over $S$ range over proper non-empty subsets of $V$.

\paragraph{Penalty Weight Justification.}
The penalty weights must be chosen to ensure that any violation of the constraints increases the energy more than the maximum possible reduction in the objective function. Let $W_{\max} = \max_{(i,j) \notin E} 1$ be the largest coefficient in the objective. The maximum objective gain achievable by violating a single degree constraint is bounded by the graph degree, which is at most $N-1$. Therefore, choosing $\lambda > W_{\max} \cdot (N-1)$ guarantees that degree violations are energetically unfavorable. For subtour elimination, the penalty $\mu$ must dominate the degree penalty to prioritize global connectivity, so we require $\mu > \lambda \cdot \max_S (|S|-1)$. These general bounds ensure a clear energy gap between feasible and infeasible states.

\paragraph{Correctness Argument.}
Minimizing $H(x,s)$ over $\{0,1\}^{n+m}$ recovers the optimal Hamiltonian cycle because (i) any assignment satisfying both the degree and subtour constraints incurs zero penalty, reducing $H$ exactly to the original objective function; (ii) any assignment violating at least one constraint incurs a penalty strictly greater than the maximum possible improvement in the objective function, due to the selection of $\lambda$ and $\mu$ as derived above. Consequently, the global minimum of $H$ must correspond to a feasible assignment that minimizes the original objective, i.e., a Hamiltonian cycle with the minimum number of missing edges. If a Hamiltonian cycle exists, the minimum energy equals zero.

\paragraph{Illustrative Example.}
Consider a graph with $N=4$ vertices $V=\{0,1,2,3\}$ and all possible edges present (the complete graph $K_4$). The number of edge variables is $N(N-1)/2 = 6$, indexed $x_0,\dots,x_5$. We consider four proper subsets $S_1=\{0,1\}$, $S_2=\{0,2\}$, $S_3=\{1,2\}$, $S_4=\{0,1,2\}$, each requiring a binary slack variable $s_1,\dots,s_4$. To match the 14-variable structure in the specification, we introduce a second slack bit per subset, yielding variables $s_{S,1}, s_{S,2}$ for $S \in \{S_1,\dots,S_4\}$, indexed $x_6,\dots,x_{13}$. 

Instantiating the general formulas with $N=4$, $\lambda=100$, and $\mu=1000$, the objective terms for non-edges (here assumed empty for a complete graph, but the formulation retains them for generality) yield coefficients of $1$ for $x_0,\dots,x_5$. The degree penalty expands to linear terms $-300 x_k$ and quadratic terms $200 x_k x_l$ for adjacent edges sharing a vertex. The subtour penalty introduces cross-terms between edge variables and slack bits, e.g., $Q_{0,6}=400$, $Q_{0,7}=800$, $Q_{6,7}=800$, along with linear slack terms $-300 s_{S,1}$ and $-600 s_{S,2}$. The constant offset evaluates to $c = 4(100)(4) + 1000 \sum_{m=1}^4 (|S_m|-1)^2 = 1600 + 1000(1+1+1+4) = 6600$. 

Solving this $14$-bit QUBO yields the optimal binary vector corresponding to a valid $4$-cycle (e.g., $x_0=x_1=x_2=x_3=1$, all others $0$) with slack variables set to zero, achieving the minimum energy. The recovered cycle visits every vertex exactly once, confirming the formulation's correctness on this instance.

\end{document}